% Source: Weinberg GR, physical PDF pp. 173--174
%         (printed pp. 148--149).
% Coverage: Section 6.10, Geodesic Deviation; equation (6.10.1).
% Figures/tables/footnotes: optional-reading footnote; no figures or tables.
% Status: source-reviewed and compile-clean.
% Uncertainties: none.

\subsection{Geodesic Deviation\optionalreading}\label{sec:6.10}

\begin{quote}
\small\emph{This section lies somewhat out of the book's main line of
development and may be omitted in a first reading.}
\end{quote}

The introduction of the curvature tensor was motivated here by the need to
construct suitable field equations for the gravitational field. However, the
curvature tensor is also useful in expressing the effects of gravitation on
physical systems.

For instance, consider a pair of nearby freely falling particles that travel
on trajectories \(x^\mu(\tau)\) and
\(x^\mu(\tau)+\xi^\mu(\tau)\), where
\(\xi^\mu\equiv\delta x^\mu\). The equations of motion are
\[
  0
  =
  \frac{\dd^2x^\mu}{\dd\tau^2}
  +\Gamma^\mu{}_{\nu\lambda}(x)
   \frac{\dd x^\nu}{\dd\tau}
   \frac{\dd x^\lambda}{\dd\tau},
\]
\[
  0
  =
  \frac{\dd^2}{\dd\tau^2}(x^\mu+\xi^\mu)
  +\Gamma^\mu{}_{\nu\lambda}(x+\xi)
   \frac{\dd}{\dd\tau}(x^\nu+\xi^\nu)
   \frac{\dd}{\dd\tau}(x^\lambda+\xi^\lambda).
\]
Evaluating the difference between these equations to first order in
\(\xi^\mu\) gives
\[
  0
  =
  \frac{\dd^2\xi^\mu}{\dd\tau^2}
  +\partial_\rho\Gamma^\mu{}_{\nu\lambda}\,
   \xi^\rho
   \frac{\dd x^\nu}{\dd\tau}
   \frac{\dd x^\lambda}{\dd\tau}
  +2\Gamma^\mu{}_{\nu\lambda}
   \frac{\dd x^\nu}{\dd\tau}
   \frac{\dd\xi^\lambda}{\dd\tau},
\]
or, in terms of covariant derivatives along the curve \(x^\mu(\tau)\) (see
Section~\ref{sec:4.9}),
\begin{equation}
  \frac{D^2\xi^\lambda}{D\tau^2}
  =
  R^\lambda{}_{\nu\rho\mu}\,
  \xi^\mu
  \frac{\dd x^\nu}{\dd\tau}
  \frac{\dd x^\rho}{\dd\tau}.
  \tag{6.10.1}\label{eq:6.10.1}
\end{equation}
Although a freely falling particle appears to be at rest in a coordinate
frame falling with the particle, a pair of nearby freely falling particles
will exhibit a relative motion that can reveal the presence of a
gravitational field to an observer that falls with them. This is of course
not a violation of the Principle of Equivalence, because the effect of the
right-hand side of Eq.~\eqref{eq:6.10.1} becomes negligible when the
separation between particles is much less than the characteristic dimensions
of the field.
